\begin{enumerate}[label=\textbf{\ansnum.\arabic*}]
  \item $8+20=28$\,ms, about $35.7$\,Hz. Use the deadline method: \verb!t_next += 0.02! and sleep for \verb!t_next - perf_counter()!, which leaves $12$\,ms of sleep.
  \item (a) \texttt{time.perf\_counter()}, (b) \texttt{time.time()} (wall-clock, human-readable), (c) \texttt{time.monotonic()}. \texttt{time.time()} can jump forward or backward when the system clock is corrected, which would give negative durations or timeouts that fire too early or too late.
  \item For example: the garbage collector runs; another process uses the CPU or an interrupt occurs (the scheduler); a slow operation such as logging, printing or a disk write; a first call that loads a module or allocates memory; thermal throttling of the CPU.
  \item Usually not: the operating system and garbage collector cause jitter of milliseconds, which is too much for a hard $1$\,ms deadline. Alternatives: a compiled language such as Rust or C++ (with a real-time kernel and real-time priority), or let the motor driver hardware run the fast loop and send only setpoints from Python at a lower rate.
  \item Pure Python calculation: about $4$\,s (the GIL runs one thread at a time), no gain from threads (use processes). Waiting for the network: about $1$\,s, because the waits overlap.
  \item The queue decouples the two rates and is thread-safe. Old frames are worthless for a robot, so a small queue avoids a growing delay. When it is full, drop the oldest (or the new) frame instead of blocking the camera.
  \item A call without a timeout can wait forever (cable pulled, server down) and freezes the loop, so the robot stops reacting. With a timeout the program notices, logs and recovers or stops safely.
  \item
\begin{verbatim}
import time
t_next = time.perf_counter()
try:
    while True:
        print("tick")
        t_next += 0.5
        time.sleep(max(0.0, t_next - time.perf_counter()))
except KeyboardInterrupt:
    print("bye")
\end{verbatim}
  \item \texttt{SIGKILL} cannot be caught, so neither the handler nor \texttt{finally} runs and the motors stay in their last state. A \textbf{watchdog} in the motor controller (stop if no command arrives within a short time) or a hardware emergency stop protects the robot.
  \item[\textbf{\ansnum.10}]
\begin{verbatim}
import numpy as np, time
durations = []
t_end = time.perf_counter() + 10
t_next = time.perf_counter()
while time.perf_counter() < t_end:
    t0 = time.perf_counter()
    step()
    durations.append(time.perf_counter() - t0)
    t_next += 0.02
    time.sleep(max(0.0, t_next - time.perf_counter()))
d = np.array(durations) * 1000
print(d.mean(), d.max(), np.percentile(d, 99))
\end{verbatim}
  If the mean is small and the maximum is many times larger, rare spikes (garbage collection, scheduling) dominate; the $99$th percentile tells how often the deadline is at risk. Decide on the maximum, not the mean (Chapter~\ref{ch:linux8}).
\end{enumerate}
